% Options for packages loaded elsewhere
\PassOptionsToPackage{unicode}{hyperref}
\PassOptionsToPackage{hyphens}{url}
\documentclass[
  11pt,
]{article}
\usepackage{xcolor}
\usepackage[left=1in, right=1in, top=1in, bottom=1in]{geometry}
\usepackage{amsmath,amssymb}
\setcounter{secnumdepth}{5}
\usepackage{iftex}
\ifPDFTeX
  \usepackage[T1]{fontenc}
  \usepackage[utf8]{inputenc}
  \usepackage{textcomp} % provide euro and other symbols
\else % if luatex or xetex
  \usepackage{unicode-math} % this also loads fontspec
  \defaultfontfeatures{Scale=MatchLowercase}
  \defaultfontfeatures[\rmfamily]{Ligatures=TeX,Scale=1}
\fi
\usepackage{lmodern}
\ifPDFTeX\else
  % xetex/luatex font selection
\fi
% Use upquote if available, for straight quotes in verbatim environments
\IfFileExists{upquote.sty}{\usepackage{upquote}}{}
\IfFileExists{microtype.sty}{% use microtype if available
  \usepackage[]{microtype}
  \UseMicrotypeSet[protrusion]{basicmath} % disable protrusion for tt fonts
}{}
\makeatletter
\@ifundefined{KOMAClassName}{% if non-KOMA class
  \IfFileExists{parskip.sty}{%
    \usepackage{parskip}
  }{% else
    \setlength{\parindent}{0pt}
    \setlength{\parskip}{6pt plus 2pt minus 1pt}}
}{% if KOMA class
  \KOMAoptions{parskip=half}}
\makeatother
\usepackage{graphicx}
\makeatletter
\newsavebox\pandoc@box
\newcommand*\pandocbounded[1]{% scales image to fit in text height/width
  \sbox\pandoc@box{#1}%
  \Gscale@div\@tempa{\textheight}{\dimexpr\ht\pandoc@box+\dp\pandoc@box\relax}%
  \Gscale@div\@tempb{\linewidth}{\wd\pandoc@box}%
  \ifdim\@tempb\p@<\@tempa\p@\let\@tempa\@tempb\fi% select the smaller of both
  \ifdim\@tempa\p@<\p@\scalebox{\@tempa}{\usebox\pandoc@box}%
  \else\usebox{\pandoc@box}%
  \fi%
}
% Set default figure placement to htbp
\def\fps@figure{htbp}
\makeatother
\setlength{\emergencystretch}{3em} % prevent overfull lines
\providecommand{\tightlist}{%
  \setlength{\itemsep}{0pt}\setlength{\parskip}{0pt}}
\usepackage[]{natbib}
\bibliographystyle{plainnat}
\usepackage{float}
\usepackage{sectsty}
\usepackage{paralist}
\usepackage{fancyhdr}
\usepackage{lastpage}
\usepackage{natbib}
\bibliographystyle{plainnat}
\usepackage{setspace}
\doublespacing
\usepackage{siunitx}
\usepackage[nottoc, numbib]{tocbibind}
\usepackage{indentfirst}
\setlength{\parindent}{24pt}
\pagestyle{fancy}
\fancyhf{}
\rfoot{Page \thepage\ of \pageref{LastPage}}
\usepackage{bookmark}
\IfFileExists{xurl.sty}{\usepackage{xurl}}{} % add URL line breaks if available
\urlstyle{same}
\hypersetup{
  hidelinks,
  pdfcreator={LaTeX via pandoc}}

\author{}
\date{\vspace{-2.5em}}

\begin{document}

\raggedright

\setlength{\parindent}{24pt}

While it is well understood that females and males are essentially
genetically identical, sex determination leads to differences in gene
expression, which ultimately leads to SSD. However, the developmental
genetic mechanisms regulating SSD remain largely unknown. Bird \&
successfully selected upon sexual dimorphism of Drosophila wingsin 1972 but did not have the
developmental-genetic tools that are available today to identify the
proximate mechanisms they manipulated. Stuart \& Rice's (2018)
experiment altered SSD over the course of 250 generations but did not
specifically select on SSD or utilize family selection as I have done
with my experiment. The scale of their design is also larger than what is feasible for a PhD thesis. As such, while
previous artificial selection experiments such as Bird \& Schaffer (1972)
and Stuart \& Rice (2018) have successfully altered the sexual
dimorphism of D. melanogaster, the resulting lineages were not
subsequently genetically, developmentally, or physiologically analyzed. NEED TO UPDATE AUDET ET AL

\{section\}Methods

\{subsection\}Initial Population and Rearing Conditions 
The initial
population of flies came from the Frankino super-colony stock. The
colony was established from more than 60 females collected from a
winery in Michigan in 2008. For the past 14 years the Frankino lab
has been varying the temperature and diet of the colony to prevent
adaptation to specific lab conditions and to maintain genetic diversity.
Throughout this experiment flies were fed a high protein diet as
described in Mirth et al. (2019). Flies were maintained in 24-hour
lighting and reared and mated at 22\si{\celsius}, unless otherwise
stated. 

\subsection{Establishing the Lineages - Generation 0} 
FVW flies were allowed to ovipositand be maintained at low density and maintained until the late pupal stage (P13/14).
Pupae were sexed and 150 single female-male pairs were
transferred as pupae to individual food vials to establish single-pair families.
Following eclosion, the adult pairs were allowed 2-5 days to mate and
oviposit, after which the parents were collected and stored dry at -80 DEGREE C.
The offspring were then sexed as pupae (P13/14), and at least eight male and eight
female pupae from each family were put into single-sex vials and allowed to
eclose at 17˚C. This generated two vials for each family: one for
females and one for males. I measured
the pupal area using uManager and a semi-automated FIJI journal ..appendix the code put on DRYAD first then maybe github.. . SSD for each family was calculated as .. put in as equation.. average log female pupal case size -- average
log male pupal case size {[}µm2{]}. I selected 20 families at
random to establish the control lineage (CL). The remaining 100 families were
ranked based on SSD. The 20 families with the highest SSD were used
to establish the high SSD lineage (HL) and the 20 families with the
lowest SSD were used to establish the low SSD lineage (LL). 

The result is 120 vials, comprising male and female vials from 60 families (20 h, 20l, 20c) each with 8 individuals (8 males, 8 females), ranked by family SSD.


INSERT ESTABLISHING FIGURE HERE

\subsection{Crossing}
WHAT AM I GOING TO DO WITH THOSE 120 VIALS. DESCRIBE SELECTION FOR THE HIGH LINE

The protocol for HL and LL is identical.

INSERT CROSSING FIGURE HERE

\{subsection\}Selection Regime - Generations 1-15
\textit{This selection regime was adapted from
Bird and Schaffer's (1972).}

.. make less abstract.. I will use the high line as an example
We ranked the
20 HL by highest SSD respectively. Thus, each family
in each selection lineage was assigned a number from 1-20, 1 being the highest SSD family. Families ranked 1-5 were reciprocally crossed
with one another in single-male-female matings, to produce twenty unique
families (Figure 2A). Ranks 6-20 were reciprocally crossed with nearest
neighbors in single-male-female matings, resulting in thirty more unique
families (Figure 2B). Because single-male-female matings may not
generate offspring, we set up two replicate matings per cross, arbitrarily
assigned mating A and mating B. Thus, I
generated a total of 50 unique crosses per lineage each generation, with
each cross represented by two single-pair matings (A and B).


Parents were collected several days after crossing and stored
in 70\% ethanol at room temperature. Larvae were left to develop at
22\si{\celsius} and sexed at P13/14. Male and female pupae were then
separated into same-sex vials from each single-pair mating, and allowed
to eclose at 17\si{\celsius}. If mating A's did not
produce eight adult males and eight adult females, then I used mating B. (This
number was lowered to six females and six males at generation 11 due to
reduced fecundity in the selected lineages). The pupal cases of females
and males from one single-pair cross per family were then measured and
the SSD for that family was calculated as described above. The 50 HL and
LL families were then ranked by descending and ascending SSD
respectively, while the 50 CL families were 'ranked' randomly.
Finally, I established single-pair matings among families ranked 1-20 in
each lineage, using the same regime described above, to generate the
next generation. This was repeated across 15 generations. To combat
potential inbreeding depression, selection was relaxed at generations 5,
10, and 15 by taking the families ranked 1-20 in the HL and LL lineages,
randomly re-assigning their ranking, and then crossing them as described
in Figure 2. 

In
contrast, the 20 CL families were 'ranked' randomly.

\subsection{Inbreeding Coefficient} 
Pedigree was embedded into the family labeling system so that it could be tracked. For all vials of the initial generation labels began with A00, with A serving as a placeholder and 00 representing generation zero. The full family ID for one such family would be A00-0082, where the dash improved readability and this just happens to be the 82nd cross labeled. For subsequent generations family IDs started with either C, H, or L (representing CL, HL, and LL respectively) followed by the generation number. The dash to improve readability would be a - or a /~ to differentiate replicate vials. The next four digits would be the parents' familial rankings, with the first two digits reflecting the dame's ranking and the last two reflecting the sire's. 

FIGURE ON LABELING

Excel code R package

\{subsection\}Characterizing SSD of Isogenized FVW Lineages
The Frankino lab has generated lineages from the same colony by isogenizing. Characterization of these lineages should hopefully provide an idea of the attainable range in morphospace

\{section\}Results 
Over nineteen months of intense selection, my team of
undergraduates and I collected and measured over 35,000 pupal cases. As
with any selection experiment, there were concerns of unintentionally
selecting upon traits affected by SSD like developmental timing and
fecundity using this selection regime.

\bibliography{export.bib}

\end{document}
